\cotefile{74}
\scanpage{19}

% L1–5
\[
\left\{\begin{aligned}
 &V\quad\text{groupe}\\
 &I\subset V\quad\text{partie de }V\\
 &V\xrightarrow{\varphi_\lambda}V_\lambda\quad(\lambda\in\Lambda)
 \quad\text{hom de groupes},\quad I_\lambda\subset V_\lambda\quad\text{partie de }V_\lambda\\
 &\deleted{\unsure{T|X}}\;Y\ \deleted{\unsure{\text{ens}}}\ \text{un }V\text{-ens.}
 \quad X_\lambda=\deleted{X}Y\mathbin{\wedge}_V(V_\lambda,\varphi_\lambda),\
 \text{d'où }\tau_\lambda:Y\mathbin{\unsure{-|\to}}X_\lambda\\
 &S=\deleted{\unsure{X_a|X}}\{a\}\amalg Y\amalg
 \overbrace{\coprod_{\lambda\in\Lambda}
 \underbrace{Y\times_V(V_\lambda,\varphi_\lambda)}_{X_\lambda}}^{X}
\end{aligned}\right.
\]

% L6
Structure de graphe sur $S$ :

% L7–10
\[
\left\{\begin{aligned}
 &\unsure{\{a\}|a|e}\ \text{est lié aux él. de } \deleted{\unsure{S_{\lambda,y}}}\,X\deleted{{}_{\unsure{\lambda|x}}}.\\
 &\deleted{\unsure{\text{les}}}\ y\in Y\ \text{est lié à }iy\quad(i\in I)\\
 &x\in X_\lambda\ \text{est lié à }ix\quad(i\in V_\lambda)\\
 &y\in Y\ \text{est lié à son image }\tau_\lambda(y)\in X_\lambda
\end{aligned}\right.
\]

% L11
et \unsure{celà|cela} épuise les arêtes.

% L12–14
À quelles conditions est-ce là un graphe trialitaire ? Je dis qu'il n'y a que les cas suivants

% L15
\[
 \text{a) }\ V\simeq\mathbb F_2^{\unsure{c_\nu|c_v|r}+1}\
 \deleted{\unsure{\nu\in\{2,3,4\}}}\ \text{où }c_\nu=
 \deleted{\unsure{\nu}}2^\nu\deleted{\unsure{+1}},\quad \nu\in\{0,1,2\}
\]

% L16
b) $I$ est une base de $V$

% L17
c)
% archive sheet number: 18
